\section{Gemeinsames Ergebnis und offene Herkunftsbrücke}

Die Kette ergibt ein klares Gesamtbild. Das vollständige ungerade Quellgitter kann lokale geladene Spinorfelder tragen, aber dann ist der volle gewünschte \(T(E_8)\)-Stromlift auf diesen Feldern nicht gegenseitig lokal. Der neutrale Quotient besitzt exakt die E8-Struktur, legt jedoch die geladene Operatorwirkung nicht fest. Ein gerades Nachbargitter realisiert den vollen Stromlift, entfernt dafür die ursprünglichen ungeraden Felder aus der lokalen Algebra. Geladene Materiefelder können dann nur als zusätzliche Module, Twistfelder, Wilsonlinien-Endpunkte oder Sektorintertwiner wiederkehren. Eine solche Rekonstruktion ist bekannte Feldtheoriestruktur, aber in den vorliegenden Verträgen noch keine aus der TFPT-Quelle ausgewählte Dynamik.

Der konstruktive Fortschritt ist dennoch erheblich: Der Randkandidat, seine Ladungen, sein vollständiger Glue, seine Parität, der explizite geladene Spinororbit, die kompatible Eichladung, der neutrale E8-Quotient, der Einheitsflussindex und die minimale lokale Nachbargitter-Reparatur sind exakt bestimmt. Ebenso exakt sind die drei entscheidenden Grenzen: Die freie Quelle wählt die E8-Energie nicht; die native endliche Fockalgebra besitzt nicht die erforderliche gemeinsame Zentrum-/Paritätswirkung; und ein Quotient oder eine Determinantenphase ersetzt keine geladene Feld- und Zeitabbildung.

Die noch fehlende Herkunftsbrücke muss deshalb in \textbf{einem} gemeinsamen Prozess liefern: die zehn lokalen Kanäle oder ihren kontrollierten Ersatz, die nichtlinearen Gewicht-eins-Ströme, den geladenen ungeraden Sektor, die Ladungen \(q,Y,g\), den Glue und die Kokzykel, den Zustand, die Energie \(V_c\) oder eine hergeleitete Alternative sowie dieselbe physische Zeit. Eine kompakte Eichverbindung wäre dabei eine zusätzliche, ausdrücklich zu begründende Struktur. Erst ihre tatsächliche Flussamplitude könnte aus der Index-Auswahlregel einen Hamilton- oder Korrelatorbeitrag machen. Bis dahin sind die Resultate exakte Struktur- und Ausschlusssätze innerhalb des bedingten Randmodells, keine Herleitung aus P1/P2, keine vierdimensionale Standardmodell-Realisierung und keine vollständige TFPT- oder TOE-Schließung.

Damit ist auch der nächste Herkunftstest klar und ohne weitere Hilfskonstruktion formuliert.

\finding{Entscheidende Herkunftslücke}{Gesucht ist kein weiterer passender Quotient, sondern ein einziger quellenselektierter Operator- und Zeittransfer, der neutrale Ströme, geladene ungerade Sektoren, Glue, Ladungen und Energie gemeinsam realisiert.}


\chapter{Sektorrekonstruktion, Paarantwort und Herkunftsgrenze}
% Integration: vor den Schlussfolgerungen zu Sektorrekonstruktion, Paarenergie und Herkunft.
\section{Quelle und Projektions-/Gedächtnisfrage}

Der Kernbefund von \src{30} ist eine Trennung, die für
alle späteren Ableitungen entscheidend bleibt. Im E8-Ziel sei
\(D=D_8\), \(s=(1/2)^8\), \(L=D\cup(D+s)\), und \(P\) projiziere auf den
\(D\)-Sektor, \(Q=1-P\) auf den Spinorsektor. Ein Spinorstrom ist
off-diagonal und kann trotz \(PAP=PBP=0\) einen geladenen Rückweg erzeugen:

\[
 P[A,B]P=PAQBP-PBQAP\ne0.
\]

Für den gemeinsamen Sugawara-Operator und die quadratischen Familien \(H_V\),
die jede Gitterladung erhalten, gilt dagegen

\[
 QHP=PHQ=0.
\]

Dann verschwinden Feshbach-Selbstenergie und Gedächtniskern,

\[
 \Sigma(z)=PHQ(z-QHQ)^{-1}QHP=0,\qquad
 K(t)=PHQe^{-itQHQ}QHP=0.
\]

Ein geladener Mess- oder Strompfad ist daher kein autonomer Übergang durch
den Hamiltonoperator. Der zweidimensionale Gegenzeuge \(P=\mathrm{diag}(1,0)\),
\(H=\mathrm{diag}(0,1)\), \(A=\sigma_x\), \(B=\sigma_y\) zeigt exakt:
\(P[A,B]P=2iP\), aber \(PHQ=0\). Für eine reale Quelle muss also gerade der
physische Projektor und sein Hamiltoncorner \(QHP\) bestimmt werden; eine
gemeinsame abstrakte Algebra genügt nicht.

Der Coxeteroperator verschiebt im Grad eins 56 integer D8-Wurzeln in den
Spinorraum. Ein D8-Readout, das unter dem tatsächlichen Clock vollständig
invariant sein soll, muss dadurch alle 248 E8-Ströme enthalten. Ein Clock-
bedingter Sektorwechsel ist aber kein Zeitwechsel. Das schließt die
Sektorerweiterung nicht aus; es verhindert nur, dass Clockcovarianz,
Stromkommutator und Gedächtnis ungeprüft gleichgesetzt werden.

\section{Native Felder, Paarjets und die positive algebraische Seite}

\src{47} verwendet die bereits festgelegte

\[
 \Gamma=\mathbb Z^{10},\quad K=\operatorname{diag}(1^9,-1),\quad
 F(r)=(r-(a\cdot r)a/2,0,-a\cdot r)
\]

mit \(V_\mathrm{aux}=K+2Kmm^TK\). In diesem Feld- und Paarvergleich bleiben die Cosinus- und Massenterme ausgeschaltet. Verwendet wird die konforme Zeit dieser bedingten masselosen Randtheorie, ergänzt um den angegebenen Viertelladungsterm. Die 64 nativen Fermiongewichte werden als
\(p_i=F(r_i)+e_R\), die 60 Mediatorgewichte als \(b_A=F(s_A)+2e_R\)
eingesetzt. Exakt 480 der 2016 Paare haben das passende Produkt und

\[
 :\Psi_i\Psi_j:=\sum_AW_{A,ij}\Phi_A,\qquad WW^T=8I_{60}.
\]

Die tatsächliche Gruppenwirkung wird über D5- und D3-Wurzeln mit einem
gemeinsamen Cocycle gelöst; die Identifikation ist daher ein
Darstellungsintertwiner und keine Dimensionsgleichheit. Die positive
Aussage lautet: Der native Tensor erscheint als lokaler Vertex-Produkt-
Tensor in dieser bedingten Randtheorie.

Die vollständige Außenquadratstruktur fehlt im Produkt am selben Punkt nicht
vollständig, sondern verschiebt sich in Jets. Für den antisymmetrisierten

\[
 A_{ij}(z)=\tfrac12[\Psi_i(z)\Psi_j(0)-\Psi_j(z)\Psi_i(0)]
\]

hat \(C_0\) Rang 60, während die Koeffizienten \(C_2\) Rang 2016 erreichen;
die 1956 dunklen Richtungen erscheinen zuerst bei Gewicht fünf. Das ist ein
exakter Feldjet-Befund, kein nativer CAR/CCR-Hamiltonoperator.

Gerade die Gram-Matrix trennt beides: Das native helle Paar und der Mediator
sind orthogonale Zustände, während ihr Coincident-Vertexbild die Matrix

\[
 \begin{pmatrix}8&\sqrt8\\\sqrt8&1\end{pmatrix}
\]

mit Determinante null besitzt. Auch die adjungierten Vertexprodukte haben
operatorwertige Pole statt kanonischer Antikommutatoren. Der native
Hamiltonoperator und sein Vakuumtheorem werden durch diese Feldabbildung
nicht transportiert. Hinzu kommt, dass die ursprünglichen $(q,Y)$-Ladungen
unter dem gewählten Dressing verschoben werden; Neutralität allein liefert
weder Amplitude noch Dynamik.

\src{48} macht die Paarantwort dennoch vollständig.
Für den radial getrennten Zustandsvektor \(J_r\) gilt mit \(A=W^TW\), Rang 60,

\[
 J_r^*J_s=G(rs)=d(rs)I+o(rs)A,
\]

\[
 d(v)=\frac{v^2(2-v)}{2(1-v)^3},\qquad
 o(v)=\frac{2-v}{2(1-v)^2}.
\]

Die Zeitantwort zerfällt damit in helle und dunkle Kanäle. Die Polarabbildung ist eine Isometrie. Für $\rho=r^2$ ist der Fehler des normierten Zeitantwort-Grenzübergangs durch $5\rho$ uniform in allen reellen Zeiten beschränkt. Bei \(r\to0\) entsteht der kontrollierte Grenzoperator

\[
 H_\mathrm{pair,0}=3P_b+5P_d=5I-A/4.
\]

Mit dem vorhandenen Viertelladungsterm wird daraus

\[
 H_\mathrm{pair}=3I-Q_2/4+2P_d.
\]

Der verbundene Paaranteil ist \(2P_d\ge0\). Er ist also keine negative helle
Bindungsenergie wie im nativen Block

\[
 \begin{pmatrix}0&\sqrt8g_c\\\sqrt8g_c&\Delta\end{pmatrix},\qquad
 E_-=(\Delta-\sqrt{\Delta^2+32g_c^2})/2<0.
\]

Das Vorzeichenproblem kann nicht durch Abziehen einer nur für $N=2$
geltenden Konstante repariert werden. Die Paarantwort ist eine präzise
positive Antwort des bestehenden gaußschen Quellraums, aber keine erzeugte
elementare Vierfermion-Wechselwirkung.

\section{Vier Sektoren, gleiche Zeit und die Grundzustandsgrenze}

\src{49} liefert die sauberste bedingte
Fermionisierung. Für

\[
 \Gamma_0=\{x\in\mathbb Z^{10}:g\cdot x\ \mathrm{gerade}\},
 \quad h=n/2,\quad f=e_9,\quad c=h+f,
\]

gibt es die vier Diskriminantklassen $(0,h,f,c)$. Die Twist-Charaktere und
die spinabhängige Arf-Kombination reproduzieren die fermionischen Spuren:

\[
 Z_F[\rho]=\frac12\sum_{a,b=0,1}
 (-1)^{ab+a\rho_t+b\rho_x}Z_B[a,b].
\]

Die inverse Matrix ist die Transponierte; die relative T-Phase ist nach
Abtrennung des gemeinsamen gravitativen Faktors \(T=\operatorname{diag}(1,1,-1,1)\).
Ungerade Felder verschieben die Sektoren \(M_0\leftrightarrow M_f\) und
\(M_h\leftrightarrow M_c\). Dadurch werden nicht nur Partitionfunktionen,
sondern auch zeitgeordnete Zweipunktfunktionen mit dem richtigen
Hamiltonoperator des Zwischensektors rekonstruiert. Gerade die Intertwiner sind nötig; vier
Zahlen oder ein Projektor allein reichen nicht.

Die verwendeten geraden Operatoren, darunter die vorhandenen Cosinus-Terme mit
Impulsen \(n,z\in\Gamma_0\), erhalten alle vier Sektoren. Das beweist keine
gemeinsame Spektralform, keine Vakuumauswahl und keinen nativen Ursprung.
Es ist eine exakte Identität innerhalb des bereits angenommenen
Lorentzgitter-/Heisenbergrahmens. Die 9R/1L-Signatur wird dabei nicht als
negative Oszillatorenergie missverstanden.

Die Folgeanalyse \src{50} betrachtet anschließend
beide tatsächlichen Grundzustände des zugewiesenen Randhamiltonoperators.
Aus der eindeutigen Zerlegung in \(F(E_8)\oplus\mathbb Z e_R\oplus\mathbb Zm\)
folgt exakt

\[
 \mathcal G_s=\operatorname{span}\{|0\rangle,|F(s)\rangle\},
\]

mit voller Lücke $1/4$. Für die 64 geladenen Felder sind die Quellladungen
$q=-3,-1,+1,+3$ mit Vielfachheiten $(15,35,13,1)$; die Additions- und
Entfernungsenergien auf den beiden Grundzuständen sind vollständig angegeben.
Im $q=-1$-Block liegen sie bei $7/4+n$ beziehungsweise $5/4+n$,
unabhängig von der Grundzustandsmischung.

Der Vergleich verwendet das unveränderte native Modell bei $g_c/\Delta=1/20$ und dessen tatsächlichen Grundzustand im Sektor $N=64$. Eine gemeinsame symmetrische Glättung von Feld und Adjungiertem macht die Antwortgewichte endlich und erhält die niedrigsten Linien. Die native Antwort hat dagegen eine positive niedrige Entfernungspolstelle
und alle Additionsenergien oberhalb der zertifizierten nativen Schwelle. Eine
positive globale Zeitskalierung, die die niedrigste Entfernungspolstelle
angleicht, erzeugt daher eine Additionspolstelle unterhalb jeder nativen
Addition. Das Verhältnis der Quellschwellen $7/5=1{,}4$ liegt weit unter
dem nativen unteren Verhältnis $>8{,}4349$. Die Differenz überlebt jedes
Vakuum in der vollständigen zweidimensionalen Grundfläche, jede Population
\(t\) und jede erlaubte Kohärenz. Ein Fit integrierter Gewichte bei \(t=1/2\)
ist möglich, aber keine gemeinsame Zeit; das Quellmoment bleibt strikt
negativ, während das native erste Moment null ist.

Damit ist der konkrete Reparaturversuch „gleiche dressed fields, gleicher
zugewiesener Rand, nur anderer Grundzustand oder Zeitfaktor" in dieser
Antwortklasse widerlegt. Das ist keine allgemeine Widerlegung aller
nichtlinearen Quellenabbildungen. Es sagt vielmehr, welche Größen eine neue
Abbildung gleichzeitig ändern und herleiten müsste.

\section{Rohträger, Herkunft und der verbleibende Nachweis}

\src{34} hält die vorgelagerte Operatorfrage offen. Wenn
ein Carrier \(H_\mathrm{car}\subset\operatorname{Ran}P_\mathrm{prim}\) liegt und
sein Involutionsoperator nur die Einschränkung desselben Projektions-
Involutionsoperators ist, dann ist seine negative Komponente identisch null.
Eine nichttriviale \(\Lambda^3E_-\otimes\Lambda^2E_+\)-Yukawaform mit
Signatur $3+2$ folgt daraus nicht. Im reinen selbstdualen CAR-Fall vertauscht
die Konjugation die beiden Polarisationen, sodass dieselbe Involution keine
3+2-Signatur auf einem J-stabilen Raum liefert. Die U(5)-Stabilität der
reinen Polarisation wählt aus sich selbst ebenfalls kein ausgezeichnetes
Drei-/Zweiplan.

Die konditionale Algebra bleibt gültig: Unter den ausdrücklich gesetzten
Voraussetzungen erzwingt eine nichtverschwindende

\[
 \Lambda^3E_-\otimes\Lambda^2E_+\to\mathbb C
\]

die Dimensionen $3$ und $2$, und die Kontraktion ist dann eindeutig bis
auf Norm und Phase. Aber diese Voraussetzungen sind nicht aus dem rohen
P1/P2-Datum hergeleitet. Geometrische Reflexion und Calderón-Projektion sind
ebenfalls nicht dasselbe: Der Calderón-Symbolanteil hängt am tangentialen
Impuls, eine punktweise geometrische Spiegelung nicht.

Der nächste Forschungsnachweis ist dadurch eng bestimmt. Er muss aus der
rohen Grenze eine lokalisierte, geladene Quelle mit der erforderlichen
Kanalzahl liefern; beide Adjungierten, Cocycle und Produkte auf demselben
Zustand und derselben Zeit müssen existieren. Er muss den integeren Carry und
den Halbspinor-Sektorwechsel in einer realen Operatorabbildung reproduzieren,
die tatsächlichen (C/J)-Wirkungen samt ihrer Phasen bestimmen und klären,
ob sie erhaltene Symmetrien, kovariante Heisenbergoperatoren oder einen
zeitabhängigen Hintergrund darstellen. Falls der relevante physische
Projektor ein Gedächtnis tragen soll, ist zusätzlich \(QHP\ne0\) für genau
diesen Quellen-Hamiltonoperator zu zeigen; ein nichtverschwindender
Stromkommutator genügt nicht.

Die gegenwärtige Evidenz ist damit positiv bei den klar begrenzten
Algebra-/Limes-/Carryaussagen und negativ bei der unbegründeten Identifikation
von Hilfsrandtheorie, nativer Wechselwirkung und roher Quelle. Genau diese
Trennung ist der belastbare Anschluss für den nächsten Dokument- und
Forschungsabschnitt.

Für die Dokumentation der Anfangszustandsfrage ist zusätzlich festzuhalten,
dass kein Resultat dieses Strangs eine Präparation stillschweigend ersetzt.
Der QWZ-Carry startet im gefüllten negativen See und entwickelt ihn unitär;
die E8-Paarantwort hält ihren Vakuumfunktional fest; die vier Sektoren halten
die Restriktion desselben geraden Operators in jedem Modul fest. Ein
Sektorlabel, eine Arf-Summe oder eine normierte Antwort ist deshalb noch kein
physischer Auswahlmechanismus. Umgekehrt darf aus dem Fehlen von
\(QHP\) bei einem abstrakten Gitterhamiltonian nicht geschlossen werden, dass
geladene Felder keine Zwischenzustände besitzen: ihre Intertwiner können
zwischen Sektoren wirken, ohne dass der ungestörte Hamiltonoperator diese
Sektoren mischt. Diese Unterscheidung ist die gemeinsame Leitplanke der
Memory-, Carry- und Rekonstruktionsberichte. Der nächste Test muss daher ein
und derselbe Operatorvergleich sein: Quelle, Zustand, geladener Operator,
Adjungierter und Zeitgenerator müssen auf einem gemeinsamen Definitionsraum
angegeben werden. Erst dann lässt sich entscheiden, ob der Carry ein lokaler
Feldprozess, ein globaler Fluxtransport oder nur eine algebraische
Sektoridentität ist.


\chapter{Das Vierphasenregister: Ausblenden oder gemeinsam kodieren}
\label{ch:register}
\section{Warum diese letzte Untersuchung nötig war}
Nach den Feld-, Gitter- und Paarvergleichen blieb eine einfache Möglichkeit offen: Vielleicht erzeugt bereits das vorhandene Phasenregister die gesuchte Wechselwirkung, wenn man seine vollständige Wirkung behält. Diese Frage wurde an der unveränderten freien QWZ-Quelle untersucht. Zwei verschiedene Reduktionsregeln müssen dabei streng getrennt werden. Die erste startet mit unabhängig präpariertem Register und Materiezustand und blendet das Register aus. Die zweite korreliert Register und Teilchenzahl und behält die geladenen Übergänge. Beide sind mathematisch zulässige Vorgänge; sie stellen unterschiedliche physische Voraussetzungen dar. \src{51}\src{52}

Die ursprüngliche Rechnung liefert vier Einteilchenblöcke. Für die Vielteilchenfrage wurde ausdrücklich die Hebung
\[
 \mathcal H_R=\C^4\otimes\operatorname{Fock}(V),\qquad
 H_R=\sum_{r=0}^3\Pi_r\otimes d\Gamma(h_r),\qquad
 h_r=b+i^ra+i^{-r}a^\dagger
\]
verwendet. $d\Gamma$ bezeichnet die zahlenerhaltende freie Vielteilchenhebung eines Einteilchenoperators. Diese Wahl ist nicht dasselbe wie $\operatorname{Fock}(\C^4\otimes V)$. Das ursprüngliche Einteilchenskript entscheidet nicht von selbst zwischen diesen Vielteilchenräumen und wählt auch deren physische Präparation nicht aus.

\section{Eine Mischung freier Entwicklungen kann nichtgaußsche Antworten tragen}
Bei einer Produktpräparation $\sigma_R\otimes\rho$ und anschließendem vollständigem Registertrace entsteht exakt
\[
 \Phi_t(\rho)=\sum_r p_r U_r(t)\rho U_r(t)^\dagger,
 \qquad p_r=\bra r\sigma_R\ket r,
 \quad U_r(t)=e^{-itd\Gamma(h_r)}.
\]
Registerkohärenzen überleben diesen Trace nicht. Eine solche Mischung ist auf der vollständigen Materiealgebra genau dann eine unitäre Entwicklung, wenn alle besetzten $U_r(t)$ bis auf skalare Phasen übereinstimmen. Für eine zusammenhängende Zeitentwicklung im hier gewählten freien Fock-Lift bleibt jeder derartige deterministische Ausgang quasifrei. Ein abschließendes, nicht konditioniertes Zurücksetzen allein des Registers ändert den bereits reduzierten Materiezustand nicht. \src{51}

Trotzdem muss eine Mischung gaußscher Zustände nicht gaußsch sein. Die konkrete Originalquelle zeigt dies bereits bei einer gewählten Einteilchenpräparation: In der Entwicklung bis zur zweiten Zeitordnung lauten die Reinheit $1-t^2$ und der summierte Wick-Defekt $-t^2/2$. Das ist eine echte Abweichung von der gaußschen Faktorisierung des reduzierten Zustands. Sie ist noch kein kohärenter quartischer Hamiltonvertex.

\finding{Ein einfaches Bild}{Das Register wählt zwischen vier Filmen. Jeder einzelne Film zeigt eine freie Entwicklung. Wer nur ihre unaufgelöste Mischung sieht, kann zusätzliche Korrelationen messen. Daraus folgt noch nicht, dass innerhalb eines Films eine neue Kraft aufgetreten ist.}

Die Aussage ist ein präziser Ausschluss dieser Produkt-und-Trace-Regel. Sie schließt weder korrelierte Kodierungen noch andere, aus der Quelle herzuleitende Wechselwirkungen aus.

\section{Eine korrelierte Kodierung erhält geladene Felder und dieselbe Zeit}
Seien $z=i$, $Z\ket r=z^r\ket r$, $S\ket r=\ket{r+1\bmod4}$ und $N$ die gesamte Materieteilchenzahl. Mit gewöhnlichen Fermionoperatoren $c_j$ definiert man
\[
 F_j=S\otimes c_j,\qquad
 C_R=Z\otimes z^N,\qquad
 J_s\ket{\psi_N}=\ket{s-N\bmod4}\otimes\ket{\psi_N}.
\]
Hier bezeichnet $C_R$ die neue kombinierte Registergröße, nicht das Coxeterelement aus den früheren Kapiteln. Die $J_s$ sind Kodierungsabbildungen und nicht die Viertelclock $J$ oder die Kopplung $J$ des Quellenkettenmodells.

Die einfachen Formeln haben einen vollständigen Inhalt:
\[
 \{F_j,F_k\}=0,\qquad
 \{F_j,F_k^\dagger\}=\delta_{jk}I,\qquad
 [N,F_j]=-F_j,
\]
\[
 F_jJ_s=J_sc_j,\qquad H_RJ_s=J_sH_s,
\]
\[
 H_s=d\Gamma(b)+z^{s-N}d\Gamma(a)+z^{-s+N}d\Gamma(a)^\dagger.
\]
Die Felder besitzen die ursprüngliche ungerade Fermionparität. Die vier Bilder der $J_s$ sind orthogonal und erschöpfen den gesamten angenommenen Quellenraum. Die Hamiltonidentität überträgt die Zeitentwicklung für alle Zeiten und damit sämtliche endlichen geladenen Operatorprodukte. Das ist stärker als die Übereinstimmung einer einzelnen komprimierten Antwort. \src{52}

Bildlich bleibt ein Rad mit vier Stellungen an die Teilchenzahl gekoppelt: Ein Teilchen zu entfernen dreht das Rad eine Stellung weiter. Radstellung plus Teilchenzahl bleiben modulo vier erhalten. Diese Korrelation ändert ausdrücklich die Voraussetzung des vorherigen Trace-Satzes; sie ist kein Widerspruch zu ihm.

\section{Die ursprüngliche Phasenlinie braucht einen zusätzlichen Registerfaktor}
Die Quelle liefert tatsächlich den Operator $D=S\otimes\Gamma(u)$, wobei $u=\operatorname{diag}(z^{q_j})$, $q_j\in\{0,1\}$, den markierten Phasenbogen beschreibt und $u^4=I$ gilt. Das Symbol $q_j$ ist hier nur eine Bogenmarkierung, keine elektrische Ladung. $D$ selbst erhält die Teilchenzahl und ist fermionisch gerade.

Der naheliegende Ausdruck $D(I\otimes c_j)$ verletzt zwischen Orten auf verschiedenen Seiten des Bogens die CAR. Die vollständige kanonische Ergänzung folgt dagegen aus einer einzigen gemeinsamen unitären Konjugation:
\[
 T=\sum_r\Pi_r\otimes\Gamma(u)^r,
 \qquad
 \boxed{F_j^u=T F_jT^\dagger=D(Z^{-q_j}\otimes c_j).}
\]
Der Registerfaktor $Z^{-q_j}$ ist entscheidend. Auch der Übergang von Stellung drei auf null stimmt, weil $\Gamma(u)^4=I$ gilt. Mit $J_s^u=TJ_s$ bleiben sowohl Feld- als auch Hamiltontransport exakt:
\[
 F_j^uJ_s^u=J_s^u c_j,
 \qquad
 H_RJ_s^u=J_s^uH_s^u,
\]
\[
 H_s^u\big|_N=d\Gamma\!\left(u^{-r}h_ru^r\right)\big|_N,
 \qquad r=s-N\bmod4.
\]
Die Konjugation verschiebt die ursprüngliche markierte Randphase an das Ende des Bogens, ohne die Hoppingbeträge zu ändern oder einen neuen Hamiltonterm einzubauen.

Auch der ursprüngliche Operator $D$ selbst geht nicht verloren:
\[
 D J_s^u=J_{s+1\bmod4}^u,\qquad C_RD=zDC_R.
\]
Die konstruierten Fermionfelder wirken innerhalb einer $C_R$-Klasse; $D$ wechselt zwischen den Klassen. Daher sind diese Klassen \emph{keine Superselektionssektoren der gesamten ursprünglichen Observablealgebra}, sobald $D$ dazugehört. Nur eine Klasse zu behalten würde einen bereits vorhandenen Quelloperator verlieren. Seine Zeitantwort wird exakt durch $H_{s+1}^u-H_s^u$ übertragen. \src{52}

\section{Höhere Fermionterme und der wirkliche räumliche Test}
Auf dem gesamten Fockraum ist $H_s^u$ im Allgemeinen kein einziger teilchenzahlunabhängiger quadratischer Operator. Im exakten Sechsmodenbeispiel der Originalquelle gilt etwa
\[
 \bra{2}H_0^u\ket{0}=\frac i2,
 \qquad \bra{2,4}H_0^u\ket{0,4}=\frac12.
\]
Die Ziffern bezeichnen besetzte Orbitale; $\ket0$ in der ersten Formel ist ein Einteilchenzustand und kein Vakuum. Ein zusätzlicher Zuschauer ändert die Phase desselben Übergangs.

Die dabei entstehenden höheren Fermionterme sind präzise eine teilchenzahlabhängige Randholonomie. In jedem festen $N$-Sektor ist die Phase eine Zahl und die Entwicklung bleibt frei. Ein Slaterdeterminant bleibt ein Slaterdeterminant. Eine durch lokale Kollision erzeugte Bindung folgt daraus nicht.

Die Lokalitätsprüfung unterscheidet zwei Klassen von Eingriffen. Zahlneutrale lokale Observablen erhalten eine uniforme Schranke für ihre Antwortkommutatoren, weil ihre Entwicklung in jedem $N$-Block einer von nur vier lokalen freien Entwicklungen entspricht. Geladene Felder wechseln $N$ und tragen in dieser Darstellung eine globale Phasenlinie. Ihre Kommutatoren erhalten die gewöhnliche Zuordnung einer vollen geladenen CAR-Algebra zu jedem einzelnen Ort nicht. Ein einzelnes globales $C_R$ ist auch kein ortsaufgelöstes Gauss-Gesetz.

Für Zylinder fester Breite lässt sich ein stärkerer Grenzfall entscheiden. Sei $X$ ein lokaler neutraler Messbereich im Abstand $R$ von der markierten Verbindung. Duhamel-Identität und lokale Ausbreitungsschranke liefern für feste Beobachtungszeiten
\[
 \|\tau_t^{H_s^u}(O_X)-\tau_t^{K_0}(O_X)\|
 \le C_X(e^{v|t|}-1)e^{-aR},
\]
wobei $K_0$ die freie unverdrehte Quelle bezeichnet und die Konstanten unabhängig von Umfang und Teilchenzahl gewählt werden können. Für $R\to\infty$ verschwindet der Unterschied in Operatornorm. Dieser Satz ist ein allgemeines Argument unter den genannten Annahmen; er wird nicht aus dem kleinen Sechsmodenbeispiel extrapoliert. \src{52}

\finding{Was damit entschieden ist}{Die korrelierte Kodierung verbindet Felder, Phasenübergänge und Zeit exakt. Im bezeichneten räumlichen Grenzfall entsteht aus ihrer Randphase keine neue lokale Kraft im Inneren. Bereits vorhandene Zustandskorrelationen und Vorgänge, die den Rand erreichen oder umwinden, sind davon getrennt. Eine andere Darstellung derselben freien Quelle ersetzt den Herkunftsnachweis einer lokalen Wechselwirkung nicht.}

Die endliche Kontrolle umfasst 183 registrierte exakte Bedingungen und 21 Quellenpins; normale und optimierte Ausführung ergeben identische Zertifikate. Die allgemeinen Beweise und die Reichweite des Bulk-Grenzsatzes sind separat ausgeschrieben. Die Wahl des Vielteilchenlifts, der physischen Felder und des Zustands sowie der Anschluss an die anders konstruierten 64 Spinor-Fermionfelder bleiben offen.

